%-------------------------------------------------------------------------------
% SECTION TITLE
%-------------------------------------------------------------------------------
\cvsection{Projects}


%-------------------------------------------------------------------------------
% CONTENT
%-------------------------------------------------------------------------------
\begin{cventries}

%---------------------------------------------------------
  \cventry
    {Master's Thesis / Industry-Academia Collaboration with Wistron NeWeb Corporation (WNC)} % Arg 1
    {A CI/CD Framework for Zero Downtime Deployment in Wi-Fi Mesh Networks}      % Arg 2
    {Feb. 2025 - Jan. 2026} % Arg 4: 注意：學位論文通常不會只有幾個月(Sep 2025-Present)，我幫你改成一般碩士專案時程，如果是別的請自行修改
    {}                    % Arg 3
    {
      \begin{cvitems}
        % 只標記 GitHub Actions, USP APIs, TR-369 (核心標準)
        \item {Integrated \textbf{GitHub Actions} with \textbf{USP APIs} to automate cloud-to-edge container delivery to Root AP via \textbf{TR-369 (USP)} standards.}
        % 只標記 C-based Controller (核心開發)
        \item {Developed a \textbf{C-based Controller} bridging USP Agent (via UDS) to Extender APs (via TCP) for synchronized mesh updates.}
        % 只標記 Linux iptables (核心技術), Zero Downtime (成果)
        \item {Engineered \textbf{Linux iptables} steering to achieve zero packet loss and \textbf{zero downtime} during Blue-Green/Canary deployments.}
      \end{cvitems}
    }

%---------------------------------------------------------
  \vspace{0.4cm} 
%---------------------------------------------------------

  \cventry
    {Industry-Academia Collaboration with Iscom Online International Info. Inc.} % Arg 1
    {Kubernetes‑based Multi-Cluster Hybrid Cloud Management System}                    % Arg 2
    {Apr. 2024 - Sep. 2024} % Arg 4
    {}                      % Arg 3
    {
      \begin{cvitems}
        % 只標記 Karmada, GitOps (核心工具)
        \item {Orchestrated public cloud and on-premises clusters using \textbf{Karmada} and \textbf{GitOps} (ArgoCD/FluxCD) for automated service propagation.}
        % 只標記 Cilium, HAProxy (核心網路組件)
        \item {Implemented \textbf{Cilium Cluster Mesh} and \textbf{HAProxy} to enable global traffic steering, cross-cluster failover, and firewall security policies.}
        % 只標記 Thanos, Prometheus, Grafana (監控三巨頭)
        \item {Engineered a unified observability stack integrating \textbf{Thanos}, \textbf{Prometheus}, and \textbf{Grafana} for centralized health monitoring.}
      \end{cvitems}
    }

%---------------------------------------------------------
\end{cventries}